\chapter*{Apresentação}
\epigraph{``\textit{Any fool can write code that a computer can understand. Good programmers write code that humans can understand}''.}{Martin Fowler}

\lettrine[lines=4, lhang=0.1, lraise=0, loversize=0.2, findent=0.1em]{\textcolor{corAzulTema}{P}}{rezado} aluno, seja bem-vindo! Este material contém diversos Capítulos, organizados de forma a guiá-lo no processo de fixação dos conteúdos aprendidos em aula, por meio de exercícios, desafios e de projetos práticos aplicados em linguagens de programação. A ordem dos Capítulos obedece a um caminho lógico que será empregado pelo professor no seu processo de aprendizagem, ou seja, a ordem dos Capítulos segue a ordem cronológica dos tópicos que serão apresentados, ensinados e treinados em laboratório.

\begin{wrapfigure}{R}{0.3\textwidth}
    \centering
    \includegraphics[width=0.25\textwidth]{imagens/david}
\end{wrapfigure}

Antes de começar, eu gostaria de me apresentar. Meu nome é David Buzatto e sou Bacharel em Sistemas de Informação pela Fundação de Ensino Octávio Bastos (2007), Mestre em Ciência da Computação pela Universidade Federal de São Carlos (2010) e Doutor em Biotecnologia pela Universidade de Ribeirão Preto (2017). Tenho interesse em algoritmos, estruturas de dados, compiladores, linguagens de programação, algoritmos em bioinformática e desenvolvimento de jogos eletrônicos. Atualmente sou professor efetivo do Instituto Federal de Educação, Ciência e Tecnologia de São Paulo (IFSP), câmpus São João da Boa Vista. A melhor forma de contatar é através do email \textcolor{blue}{\href{mailto:davidbuzatto.ifsp@gmail.com}{davidbuzatto.ifsp@gmail.com}}.

Para que você possa aproveitar a leitura deste documento de forma plena, vale a pena entender alguns padrões que foram utilizados neste texto. As três caixas apresentadas abaixo serão empregadas para mostrar, a você leitor, respectivamente, boas práticas de programação, conteúdos complementares para melhorar e aprofundar seu aprendizado e, por fim, itens que precisam ser tratados com cuidado ou que podem acarretar em erros comuns de programação.

\begin{boaPratica}
    Essa é uma caixa de ``Boa Prática''.
\end{boaPratica}

\begin{saibaMais}
    Essa é uma caixa de ``Saiba Mais''.
\end{saibaMais}

\begin{atencao}
    Essa é uma caixa de ``Atenção''.
\end{atencao}

Além disso, diversos exercícios conterão caixas de entrada e/ou saída, em que serão apresentados exemplos de dados de entrada para o problema proposto e de saída que devem ser obtidos após o processamento da entrada fornecida como exemplo.

Note também que este documento foi escrito de forma quase coloquial, com o objetivo de conversar com você e não com o objetivo de ser um material de pesquisa.

É de suma importância que você resolva cada um dos exercícios básicos de cada Capítulo, visto que a utilização de uma linguagem de programação, e mais importante ainda, a obtenção de maturidade no desenvolvimento de algoritmos, é ferramenta primordial para o seu sucesso profissional e intelectual na área da Computação.

Um ponto importante sobre os exercícios é que todas as saídas apresentadas ao usuário, quando feitas em modo texto, serão feitas sem usar acentos. Ou seja, se um programa precisar exibir uma palavra acentuada, algo como, por exemplo, ``Olá'', você deverá digitar tal palavra sem usar o acento, ficando ``Ola''. O principal motivo para essa restrição é que normalmente a página de códigos utilizada para executar o programa em modo texto diverge da codificação utilizada na compilação, criando inconsistências, principalmente ao se usar as linguagens C e C++ em ambiente Windows. Por isso, nas entradas e saídas de todos os enunciados, você perceberá que faltarão acentos nas palavras acentuadas, o que é feito de propósito visto essa ``restrição''. Há como contornar tal característica, mas isso envolve alguns detalhes que fogem do escopo deste documento e que lhe forçaria a inserir código no seu programa que não faz parte da solução nem dos conceitos que devem ser aprendidos. Usando as linguagens Java ou Python não haverá tal problema, entretanto, será mantido o padrão de não usar acentos para todos os exercícios. Essas quatro linguagens de programação foram citadas pois serão linguagens que, dependendo da disciplina em que essa lista for usada, poderão estar sendo ensinadas/aplicadas. Os programas de exemplo conterão acento normalmente.

Como última observação, vale ressaltar que este documento será constantemente atualizado, sendo assim, sempre obtenha a última versão no local indicado pelo professor. Espero que este material seja útil!

\vspace*{\fill}

\FloatBarrier
\begin{figure*}[!htbp]
    \centering
    \includegraphics[scale=0.1]{imagens/logoCC}
    \\
    Este trabalho está licenciado sob uma Licença Creative Commons Atribuição-NãoComercial-SemDerivações 4.0 Internacional. Para ver uma cópia desta licença, visite \textcolor{blue}{\url{http://creativecommons.org/licenses/by-nc-nd/4.0/}}.
\end{figure*}
\FloatBarrier
